\chapter{State of the Art}
\label{chap:estado-da-arte}

\section{Introduction}
\label{chap2:sec:intro}

This chapter reviews the body of practice and tooling relevant to the
problem addressed by this project: how to automatically validate a
Jakarta EE \ac{REST} \ac{API}, build after build, against a moving
reference. It is organised around three axes. Section~\ref{sec:soa-tests}
surveys the main families of test frameworks used in the Java world,
from low-level unit testing to in-container and \ac{HTTP}-level testing.
Section~\ref{sec:soa-cicd} discusses the \ac{CI}/\ac{CD} platforms that
are commonly used to orchestrate those tests, with particular focus
on GitLab, which is the platform used by the \ac{SEGAL} laboratory and
which directly constrains the design described later in the report.
Section~\ref{sec:soa-diff} reviews differential and regression testing
approaches that influenced the design of the harnesses presented in
Chapter~\ref{chap:caso-epos}. Section~\ref{chap2:sec:concs} closes the
chapter.

\section{Test Frameworks for the Java Ecosystem}
\label{sec:soa-tests}

The Java test landscape can be roughly partitioned into four layers:
unit testing, integration testing, container testing and \ac{API}/
\ac{HTTP} testing. Each layer answers a different question about the
system under test, and modern projects typically combine several of
them.

\subsection{Unit Testing: JUnit and TestNG}
\label{sec:junit}

\textbf{JUnit 5}~\cite{junit5} is the de facto standard for unit testing
in Java. Its modular architecture, \emph{Jupiter} as the
programming/extension model, \emph{Vintage} for backward compatibility
with JUnit 4, and the \emph{Platform} as the runner, allows it to
host a wide variety of test engines and to integrate with virtually
every build tool and \ac{IDE}. The \texttt{@Test}, \texttt{@ParameterizedTest}
and \texttt{@Nested} annotations, together with a rich set of
\emph{lifecycle} and \emph{extension} hooks, cover the vast majority of
unit testing needs.

\textbf{TestNG}~\cite{testng} is the historical alternative. While its
adoption has declined in greenfield projects since the release of JUnit
5, it still has advantages in scenarios that require sophisticated test
ordering and parallel execution out of the box, and remains common in
large enterprise code bases.

For this project, JUnit 5 was chosen because of its first-class support
in Maven Surefire/Failsafe, its native integration with the Arquillian
Jupiter container and the simplicity of running it inside a GitLab
pipeline, which natively understands the JUnit XML reporting format.

\subsection{Integration Testing: Maven Failsafe and Testcontainers}
\label{sec:failsafe}

Unit tests answer the question ``does this class behave correctly in
isolation''. Integration tests answer the question ``do these classes
behave correctly when wired together with their real collaborators''.
In a Maven-based Jakarta EE project, integration tests are conventionally
named \texttt{*IT.java} and executed by the \emph{Failsafe} plugin
during the \texttt{verify} phase, which is distinct from the
\texttt{test} phase used by Surefire for unit tests.

\textbf{Testcontainers}~\cite{testcontainers} is the most widely adopted
solution for spinning up real infrastructure, databases, message
brokers, browsers, application servers, inside Docker containers
purely for the duration of a test run. It is increasingly used as a
replacement for in-memory fakes (such as H2 acting as PostgreSQL),
because the parity between test and production environments is much
higher. Testcontainers was considered for the prototypes described in
Chapter~\ref{chap:metodo-arq} but ultimately not adopted, because the
production database fixture is hosted inside the laboratory's network
and is already reachable from the GitLab Runner via the \ac{VPN}; a
second, containerised copy would only complicate the picture.

\subsection{Container Testing: Arquillian}
\label{sec:arquillian}

\textbf{Arquillian}~\cite{arquillian} is a container testing framework
that started in the Java EE world and has since followed the rename to
Jakarta EE. Its key idea is the \emph{micro-deployment}: each test
method ships its own small \ac{WAR}/\ac{JAR}, packaged in code through
\emph{ShrinkWrap}, that is deployed into a real container (managed,
embedded or remote) just before the test runs and undeployed
immediately after. This allows \ac{JAX-RS} resources, \ac{CDI} beans
and \ac{JPA} entities to be tested with the same lifecycle and the
same configuration as in production, while avoiding the cost of
maintaining a fully populated server.

Arquillian was used in the first prototype (Section~\ref{sec:javaeeapp})
to validate the \ac{CI}/\ac{CD} pipeline, but was deliberately
\emph{not} adopted for the production \ac{API}. Two reasons motivated
this choice: first, the production application is far too large to
benefit from per-test micro-deployments; and second, the laboratory's
real need was not to test a single deployment in isolation, but to
compare \emph{two} deployments of slightly different code at the same
time, which is more naturally expressed by spawning two Payara Micro
processes in parallel.

\subsection{\ac{API}-level Testing: REST Assured, Postman/Newman}
\label{sec:apitests}

\textbf{REST Assured}~\cite{restassured} is a Java \ac{DSL} for testing
\ac{REST} services, with a fluent
``\texttt{given().\textellipsis when().\textellipsis then()}'' syntax
heavily inspired by behaviour-driven development. It is the most
common choice for \ac{HTTP}-level tests written in the same language
as the system under test, and it integrates with JUnit 5 transparently.
It was considered as an alternative to writing a custom \ac{HTTP}
harness with \texttt{java.net.http.HttpClient}, but in the end the
\emph{differential} nature of the tests, the same request, sent to
two endpoints, with two responses to be diffed, mapped more
naturally onto plain \texttt{HttpClient} calls.

\textbf{Postman}~\cite{postman} is a graphical \ac{API} client widely
used by developers and \ac{QA} teams to manually exercise endpoints.
Its companion command-line runner \textbf{Newman}~\cite{newman} can
replay Postman \emph{collections} (JSON files describing a set of
requests and assertions) inside a \ac{CI} pipeline. This is an
attractive solution for teams in which the \ac{API} clients are
maintained by non-developers, but its assertion model
(JavaScript-based, embedded in the collection) makes diffing two
responses against each other awkward.

\subsection{Differential Testing}
\label{sec:soa-diff}

A specific family of testing approaches that directly inspired the
work described later in this report is \emph{differential testing}.
Rather than asserting that the system under test produces a specific
expected response, differential testing relies on the existence of a
``reference'' implementation (or version) and asserts only that the
two implementations agree on a shared set of inputs.

The canonical published example in the service world is
\textbf{Diffy}~\cite{diffyTwitter}, a proxy released by Twitter
Engineering that sits in front of a candidate service and a primary
service and replays each incoming request against both, reporting any
divergence in the responses. Diffy and its descendants are designed
for the case in which a new version of a service must be validated by
comparing its behaviour against the version it is about to replace,
exactly the case faced by the \ac{SEGAL} laboratory during the
refactor of its \ac{API}.

A second family of approaches is so-called \emph{snapshot testing},
popularised in the JavaScript world by Jest and now common in other
ecosystems. Snapshot testing serialises the output of a piece of code
to a file the first time the test runs, and on every subsequent run
asserts that the new output matches the stored snapshot. It scales
poorly when the output of the system depends on a live database or
when the legitimate response can vary slightly between runs, both of
which are true for the \ac{EPOS} \ac{API}.

The approach taken in this project is closer to Diffy in spirit than
to snapshot testing: rather than freezing a single reference, the
reference is always the parent commit of the build currently under
test, which keeps the comparison meaningful across long refactor
branches without requiring a human to update fixtures.

\section{Continuous Integration and Continuous Delivery}
\label{sec:soa-cicd}

\subsection{Foundations}

The notion of \ac{CI}/\ac{CD} as articulated by Humble and
Farley~\cite{humbleFarley2010continuous} is the practice of integrating
code into a shared mainline frequently, typically several times a
day, and of keeping that mainline in a state from which a deployable
artifact can be produced at any time, through a fully automated
pipeline. The pipeline is canonically organised into a small number of
stages, build, test, deploy, with strict gating between them:
no stage runs unless the previous one has succeeded, and a failure
stops the pipeline and surfaces a clear signal to the team.

The same publication popularised the idea of the
\emph{deployment pipeline} as the single source of truth about a
build's quality. In a mature setup, the artifact produced by the
\texttt{build} stage is the exact same artifact that flows through the
test stages and eventually to deployment, ensuring there is no drift
between what was tested and what was shipped. The same principle is
applied in this project: the \ac{WAR} built once in the first stage of
the pipeline is consumed verbatim by the testing stages later on.

\subsection{Platforms}

Three platforms dominate the modern \ac{CI}/\ac{CD} space and are
worth comparing in the context of the laboratory's choice.

\textbf{Jenkins}~\cite{jenkins} is the long-standing, plugin-based,
self-hosted automation server. Its strengths are flexibility and
breadth of plugin coverage; its main drawbacks are that pipelines are
historically described in Groovy (with the related learning curve) and
that the maintenance burden of the master/agent topology is
non-trivial.

\textbf{GitHub Actions}~\cite{githubActions} is a hosted \ac{CI}/\ac{CD}
service that integrates tightly with repositories hosted on GitHub.
Its pipelines are described in \ac{YAML} and its ecosystem of reusable
\emph{actions} is large. It is a very strong default for open source
projects, but it cannot be self-hosted as an end-to-end platform; the
control plane runs on GitHub's infrastructure.

\textbf{GitLab \ac{CI}/\ac{CD}}~\cite{gitlabci} is the \ac{CI} system
integrated into GitLab. Pipelines are described in a single
\texttt{.gitlab-ci.yml} file at the root of the repository; runners are
separate processes that pick up jobs from the GitLab server, and can
be \emph{shared} (provided by gitlab.com) or
\emph{specific}~\cite{gitlabRunner}, registered against a particular
project or group. The feature that matters most for this project is that GitLab
itself can be \emph{self-managed}~\cite{gitlabSelfManaged}: the entire
control plane can be installed inside the laboratory's network, and
combined with one or more locally-hosted runners can be used to
build, test and deploy applications that depend on internal resources
such as the laboratory's \ac{GNSS} database, without exposing
those resources to the public Internet.

That last property is the deciding factor. The \ac{SEGAL} laboratory
operates a self-managed GitLab instance and a GitLab Runner host
provisioned specifically for the \ac{EPOS} \ac{API}; that runner is
inside the laboratory's network and can reach the production database
directly. The choice of GitLab \ac{CI}/\ac{CD} as the orchestration
platform for this project is therefore not an abstract preference but
a consequence of the operating environment.

\section{Conclusion}
\label{chap2:sec:concs}

The state of the practice in Java \ac{API} testing offers a rich
catalogue of tools at each layer of the stack, and a small number of
mature \ac{CI}/\ac{CD} platforms in which to orchestrate them. From
this catalogue, the project adopts JUnit 5 as the unit/integration
test framework, plain \texttt{java.net.http.HttpClient} as the
\ac{HTTP} client for the differential harness, and GitLab \ac{CI} as
the orchestration platform, the last choice being effectively
mandated by the laboratory's infrastructure. Differential testing, as
embodied by tools like Diffy, provides the conceptual model on top of
which the bespoke harnesses described in Chapter~\ref{chap:caso-epos}
are built. The next chapter details the concrete stack of technologies
on which the implementation rests.
